%% notes-tex/database/044-organism-similarity/044-organism-similarity.tex
%% [[experiments.044-organism-similarity]]
%% https://github.com/Mjvolk3/torchcell/tree/main/notes-tex/database/044-organism-similarity/044-organism-similarity.tex
%% SOURCE: experiments/044-organism-similarity/scripts/08_notes_tex_numbers.py (numbers.tex, figures.tex, tables/)

\documentclass[10pt,a4paper]{article}
\usepackage[draft]{tcdoc}
\usepackage{float}

\newcolumntype{L}[1]{>{\raggedright\arraybackslash}p{#1}}

%% GENERATED FILE -- do not hand-edit.
%% SOURCE: experiments/044-organism-similarity/scripts/08_notes_tex_numbers.py
\newcommand{\nOrganisms}{34}
\newcommand{\nProteomes}{33}
\newcommand{\nPairs}{528}
\newcommand{\nIsolates}{1,011}
\newcommand{\nOrfs}{6,005}
\newcommand{\isoMedMedId}{99.72}
\newcommand{\isoMedMedDiv}{0.28}
\newcommand{\isoMedMeanId}{98.04}
\newcommand{\isoMinMeanId}{95.55}
\newcommand{\isoMaxMeanId}{99.94}
\newcommand{\isoEqLenId}{99.58}
\newcommand{\litIdLo}{99.90}
\newcommand{\litIdHi}{99.91}
\newcommand{\litDivMid}{0.09}
\newcommand{\litSitesLo}{10,000}
\newcommand{\litSitesHi}{12,000}
\newcommand{\litResiduesM}{11.7}
\newcommand{\aaiHumChimp}{98.14}
\newcommand{\aaiShortHumChimp}{98.1}
\newcommand{\aaiMedHumChimp}{99.20}
\newcommand{\rbhHumChimp}{18,461}
\newcommand{\ofHumChimp}{0.920}
\newcommand{\scHumChimp}{98.07}
\newcommand{\aaiHumMacaque}{95.29}
\newcommand{\aaiShortHumMacaque}{95.3}
\newcommand{\aaiMedHumMacaque}{96.90}
\newcommand{\rbhHumMacaque}{18,064}
\newcommand{\ofHumMacaque}{0.900}
\newcommand{\scHumMacaque}{94.85}
\newcommand{\aaiHumMouse}{83.70}
\newcommand{\aaiShortHumMouse}{83.7}
\newcommand{\aaiMedHumMouse}{86.50}
\newcommand{\rbhHumMouse}{16,925}
\newcommand{\ofHumMouse}{0.843}
\newcommand{\scHumMouse}{80.92}
\newcommand{\aaiHumChicken}{72.12}
\newcommand{\aaiShortHumChicken}{72.1}
\newcommand{\aaiMedHumChicken}{73.60}
\newcommand{\rbhHumChicken}{13,170}
\newcommand{\ofHumChicken}{0.731}
\newcommand{\scHumChicken}{68.16}
\newcommand{\aaiScerSpar}{92.15}
\newcommand{\aaiShortScerSpar}{92.1}
\newcommand{\aaiMedScerSpar}{93.20}
\newcommand{\rbhScerSpar}{5,367}
\newcommand{\ofScerSpar}{0.969}
\newcommand{\scScerSpar}{92.09}
\newcommand{\aaiScerSeub}{83.55}
\newcommand{\aaiShortScerSeub}{83.5}
\newcommand{\aaiMedScerSeub}{84.70}
\newcommand{\rbhScerSeub}{5,136}
\newcommand{\ofScerSeub}{0.955}
\newcommand{\scScerSeub}{83.64}
\newcommand{\aaiSdCbs}{99.56}
\newcommand{\aaiShortSdCbs}{99.6}
\newcommand{\aaiMedSdCbs}{99.80}
\newcommand{\rbhSdCbs}{5,037}
\newcommand{\ofSdCbs}{0.980}
\newcommand{\scSdCbs}{99.59}
\newcommand{\aaiScerEcoli}{36.84}
\newcommand{\rbhScerEcoli}{391}
\newcommand{\ofScerEcoli}{0.091}
\newcommand{\nGenus}{3}
\newcommand{\aaiMedianGenus}{84.2}
\newcommand{\aaiMinGenus}{83.5}
\newcommand{\aaiMaxGenus}{92.1}
\newcommand{\nMammals}{28}
\newcommand{\aaiMedianMammals}{84.3}
\newcommand{\aaiMinMammals}{83.5}
\newcommand{\aaiMaxMammals}{98.1}
\newcommand{\nMammalBird}{8}
\newcommand{\aaiMedianMammalBird}{72.1}
\newcommand{\aaiMinMammalBird}{71.5}
\newcommand{\aaiMaxMammalBird}{72.2}
\newcommand{\nYcotina}{75}
\newcommand{\aaiMedianYcotina}{47.4}
\newcommand{\aaiMinYcotina}{46.0}
\newcommand{\aaiMaxYcotina}{99.6}
\newcommand{\nBacteria}{45}
\newcommand{\aaiMedianBacteria}{40.5}
\newcommand{\aaiMinBacteria}{38.8}
\newcommand{\aaiMaxBacteria}{72.7}
\newcommand{\nYeastBacteria}{130}
\newcommand{\aaiMedianYeastBacteria}{36.5}
\newcommand{\aaiMinYeastBacteria}{35.1}
\newcommand{\aaiMaxYeastBacteria}{38.2}
\newcommand{\scN}{253}
\newcommand{\scR}{0.998}
\newcommand{\scDiff}{1.66}
\newcommand{\sdAai}{47.1}
\newcommand{\sdOf}{0.595}
\newcommand{\sdOfPct}{59.5}
\newcommand{\sdRbh}{3,072}
\newcommand{\sdGenes}{5,159}
\newcommand{\cbsAai}{47.1}
\newcommand{\minVertAai}{71.5}
\newcommand{\nVertPairs}{36}
\newcommand{\ycoAaiMin}{46.4}
\newcommand{\ycoAaiMax}{59.2}
\newcommand{\ycoOfMin}{0.410}
\newcommand{\ycoOfMax}{0.868}
\newcommand{\animAaiMin}{41.0}
\newcommand{\animAaiMax}{41.4}
\newcommand{\bacAaiMin}{35.1}
\newcommand{\bacAaiMax}{37.0}
\newcommand{\bacOfMax}{0.111}
\newcommand{\nSccPartners}{32}
\newcommand{\nYco}{10}
\newcommand{\nAnimal}{9}
\newcommand{\nBac}{10}
\newcommand{\nYard}{8}
\newcommand{\nLeaves}{30}
\newcommand{\nNamedClades}{29}
\newcommand{\nNoClade}{49}
\newcommand{\nAllPairs}{510,555}
\newcommand{\medAllPairs}{1.02}
\newcommand{\nBelow}{27}
\newcommand{\nNamedWithPairs}{28}
\newcommand{\cladeMedMin}{0.0015}
\newcommand{\cladeMedMinName}{22. Far East Russian}
\newcommand{\cladeMedMax}{1.03}
\newcommand{\cladeMedMaxName}{M3. Mosaic region 3}
\newcommand{\medCladeMed}{0.44}
\newcommand{\pooledWithin}{0.21}
\newcommand{\minBoxPairs}{10}
\newcommand{\nBoxClades}{24}
\newcommand{\nStripClades}{5}
\newcommand{\noCladeMed}{1.15}
\newcommand{\nearestStrain}{ADT}
\newcommand{\nearestDiv}{0.0005}
\newcommand{\nearestClade}{M3. Mosaic region 3}
\newcommand{\nearestCladeN}{113}
\newcommand{\nNegBranches}{0}
\newcommand{\nClsGenus}{3}
\newcommand{\nClsMammals}{28}
\newcommand{\nClsMammalBird}{8}
\newcommand{\nClsYcotina}{75}
\newcommand{\nClsBacteria}{45}
\newcommand{\nClsYeastBacteria}{130}
\newcommand{\nClsYeastAnimal}{117}
\newcommand{\nClsAnimalBacteria}{90}
\newcommand{\nClsBasidioYeast}{13}
\newcommand{\nClsBasidioAnimal}{9}
\newcommand{\nClsBasidioBacteria}{10}
\newcommand{\nClasses}{11}

%% GENERATED FILE -- do not hand-edit.
%% SOURCE: experiments/044-organism-similarity/scripts/08_notes_tex_numbers.py
\newcommand{\figTransfer}{figures/fig6_main_transfer_ladder_2026-10-19-20-33-31.pdf}
\newcommand{\figClaim}{figures/fig1_claim_protein_divergence_2026-10-19-20-33-31.pdf}
\newcommand{\figHeatmaps}{figures/fig2_cladogram_heatmaps_2026-10-19-20-33-31.pdf}
\newcommand{\figTangle}{figures/fig3_taxonomy_vs_aai_nj_2026-10-19-20-33-31.pdf}
\newcommand{\figClades}{figures/fig4_clade_tree_within_distance_2026-10-19-20-33-31.pdf}
\newcommand{\figOrthologs}{figures/fig5_ortholog_fraction_vs_aai_2026-10-19-20-33-31.pdf}


\title{\vspace{-1.5em}Organism Similarity: Is There More Genetic Difference Within Yeast
Than Within Mammals?}
\author{Michael Volk}
\date{\today}

\begin{document}
\maketitle
\thispagestyle{empty}

\noindent
Yes, at the genus and the subphylum, which is where the claim is made (Peris et al. 2023,
Fig.~3c), and the data here reproduce it (terms in Table~\ref{tab:vocab}). Within
\org{Saccharomyces}, S288C/\org{S. paradoxus} AAI is \aaiShortScerSpar\% and
S288C/\org{S. eubayanus} \aaiShortScerSeub\%, so the genus spans the distance from
human/macaque (\aaiShortHumMacaque\%) to human/mouse (\aaiShortHumMouse\%). The median of
the \nYcotina{} Saccharomycotina pairs outside the genus is \aaiMedianYcotina\% AAI, below
every one of the \nVertPairs{} vertebrate pairs (lowest \minVertAai\%). Within
\org{S.~cerevisiae}, an isolate's typical ORF is \isoMedMedId\% identical to S288C at the
protein level (median over the \nIsolates{} isolates of each isolate's per-ORF median), and
\isoMedMeanId\% when every frameshifted or indel-bearing ORF is charged in full; one human
genome differs from GRCh38 at \litIdLo{} to \litIdHi\% protein identity, a literature value
not computed here. Measured: AAI and ortholog fraction over reciprocal best mmseqs2 hits for
the \nPairs{} pairs of \nProteomes{} per-gene proteomes, EC-4 repertoires from one
eggNOG-mapper run, isolate protein identity from Peter et al. 2018's reference-inferred
alleles, and the clade structure of Peter's SNP distance matrix. For carrying a model
trained on the isolates to \org{P. kudriavzevii} SD108, the numbers that matter are
Figure~\ref{fig:transfer}'s: SD108 is \sdAai\% AAI from S288C, beyond every vertebrate
pair, and \sdOfPct\% of its genes have a one-to-one ortholog in S288C.

\vspace{0.5em}
\tccontents
\clearpage

%% notes-tex/database/044-organism-similarity/sections/vocabulary.tex
%% [[experiments.044-organism-similarity]]
%% https://github.com/Mjvolk3/torchcell/tree/main/notes-tex/database/044-organism-similarity/sections/vocabulary.tex
%% SOURCE: definitions as implemented in experiments/044-organism-similarity/scripts/03_mmseqs_rbh.py, 04_emapper.py, 05_trees.py, 06_isolate_protein_identity.py

\section{Controlled vocabulary\secstatus{todo}}\label{sec:vocab}

\begin{table}[H]\centering\footnotesize
\caption[]{Terms used in the figures and tables, in alphabetical order.}
\label{tab:vocab}
\begin{tabular}{@{}L{38mm} L{140mm}@{}}
\toprule
Term & Definition \\
\midrule
AAI & Average amino-acid identity: the mean percent identity over the reciprocal best hits
of two per-gene proteomes (mmseqs2 local alignments, \texttt{03\_mmseqs\_rbh.py}). \\
Amniote & In this organism set, a mammal or the chicken. \\
EC-4 Jaccard & Jaccard index of the sets of complete, four-level EC numbers that one
eggNOG-mapper run assigns to the two proteomes (\texttt{04\_emapper.py}). \\
Ortholog fraction & Reciprocal best hits over the gene count of the smaller per-gene
proteome, $n_\mathrm{RBH}/\min(n_i, n_j)$: the share of the smaller organism's genes that
have a one-to-one ortholog in the other. \\
Ortholog Jaccard & Reciprocal best hits over the union of the two per-gene proteomes,
$n_\mathrm{RBH}/(n_i + n_j - n_\mathrm{RBH})$. \\
Per-gene proteome & One protein per gene, the longest, keyed by NCBI GeneID where the
annotation carries it (\texttt{02\_proteome\_prep.py}). \\
Peter clade & The clade an isolate carries in Peter et al. 2018's Supplementary Table 1,
with the Wine/European subclades merged; isolates the table gives no clade form one more
group. \\
Reciprocal best hit (RBH) & A pair of proteins, one from each proteome, each the other's
best mmseqs2 hit; the proxy for a one-to-one ortholog. \\
Single-copy OG proxy & AAI restricted to Eukaryota-level eggNOG orthologous groups that
hold exactly one protein in each member of the pair; stands in for Peris et al. 2023's
BUSCO-restricted AAI. \\
SNP distance & Percent non-identical bases between two isolates in Peter et al. 2018's
released distance matrix, heterozygous differences half-weighted. \\
\bottomrule
\end{tabular}
\end{table}

%% notes-tex/database/044-organism-similarity/sections/figures.tex
%% [[experiments.044-organism-similarity]]
%% https://github.com/Mjvolk3/torchcell/tree/main/notes-tex/database/044-organism-similarity/sections/figures.tex
%% SOURCE: experiments/044-organism-similarity/scripts/07_figures.py (panels), 08_notes_tex_numbers.py (every number)

\clearpage
\section{Figures\secstatus{todo}}\label{sec:figures}

\src{07\_figures.py fig6\_main\_transfer\_ladder; s288c\_centered.csv}
\tcfig{\figTransfer}{%
\textbf{The training isolates reach human/chimp protein divergence from S288C;
\org{P. kudriavzevii} SD108 lies beyond every vertebrate pair, with \sdOfPct\% of its
genes in one-to-one orthology with \org{S. cerevisiae}.}
(a) Protein divergence from S288C (100 minus identity, log scale, inverted; identity on the
right), columns ordered by distance, colored by the partner's taxonomic group (legend).
SD108 (large point, black edge) is \sdAai\% AAI from S288C, below the lowest of the
\nVertPairs{} vertebrate pairs (\minVertAai\%; human/chicken \aaiShortHumChicken\%). The
\nIsolates{} isolates appear twice: circles are each isolate's median over ORFs of protein
identity to S288C (median over isolates \isoMedMedId\%), triangles its mean over ORFs, which
charges every gap column of a frameshifted or indel-bearing ORF (median \isoMedMeanId\%,
lowest \isoMinMeanId\%, against human/chimp AAI \aaiHumChimp\%). The other \nYco{}
Saccharomycotina hosts sit at \ycoAaiMin{} to \ycoAaiMax\% AAI, the \nAnimal{} mammals and
chicken at \animAaiMin{} to \animAaiMax\%, the \nBac{} bacteria at \bacAaiMin{} to
\bacAaiMax\%. Right of the dashed line, the same axis measured from human over \nYard{}
partners, as a yardstick (mammals blue, chicken gray). Gray band: one human genome vs GRCh38
at peptide-altering sites (1000 Genomes Project 2015; literature value).
(b) Ortholog fraction with S288C for the organisms of (a) without the isolates, same order
and colors. SD108 shares \sdRbh{} reciprocal best hits with S288C, \sdOfPct\% of its
\sdGenes{} genes; the other Saccharomycotina range from \ycoOfMin{} to \ycoOfMax{}, and no
bacterium exceeds \bacOfMax{}.}{fig:transfer}

\clearpage
\src{07\_figures.py fig1\_claim\_protein\_divergence}
\tcfig{\figClaim}{%
\textbf{\org{Saccharomyces} reaches the human/mouse distance, and the median
Saccharomycotina pair lies beyond human/chicken.}
Each point is a pair; the y axis is 100 minus protein identity (log scale, inverted;
identity on the right). Within \org{Saccharomyces} AAI runs from \aaiMinGenus{} to
\aaiMaxGenus\% (\nGenus{} pairs), within mammals from \aaiMinMammals{} to \aaiMaxMammals\%
(\nMammals{} pairs), mammal vs chicken \aaiMinMammalBird{} to \aaiMaxMammalBird\%
(\nMammalBird{}); the \nYcotina{} Saccharomycotina pairs outside the genus have a median of
\aaiMedianYcotina\%, the \nBacteria{} bacterial pairs \aaiMedianBacteria\% and the
\nYeastBacteria{} Saccharomycotina--bacteria pairs \aaiMedianYeastBacteria\%. Within
\org{S. cerevisiae} (\nIsolates{} isolates), each point is the mean over S288C ORFs of an
isolate's protein identity to S288C with every gap column charged; the per-isolate median
over ORFs gives a median divergence of \isoMedMedDiv\%. Every other column is AAI over
reciprocal best hits. Open diamonds: the single-copy OG proxy. Point face and edge colors
are the taxonomic groups of the two members, so a within-group pair is one solid color and
a cross-group pair carries both. Gray band: one human genome differs from GRCh38 at
\litSitesLo{} to \litSitesHi{} peptide-altering sites (1000 Genomes Project 2015; literature
value), divided by the \litResiduesM{} million residues of the human per-gene proteome.}{fig:claim}

\clearpage
\src{07\_figures.py fig2\_cladogram\_heatmaps}
\tcfig{\figHeatmaps}{%
\textbf{The \nOrganisms{} organisms on three similarity axes, in NCBI taxonomy order.}
Left: NCBI taxonomy cladogram (topology only; ete3 NCBITaxa, \texttt{05\_trees.py}); the
color strip is each organism's first group in the experiment plan (legend: reference,
established host, emerging host, CABBI, \org{Saccharomyces}, mammal, outgroup). Heatmaps,
rows and columns in cladogram order: AAI, ortholog Jaccard and EC-4 Jaccard (eggNOG-mapper
on eggNOG 5.0.2, \texttt{04\_emapper.py}); color scales above each map. Gray: no proteome
(\org{A. limacinum} has no annotated assembly at NCBI) or the diagonal.}{fig:heatmaps}

\clearpage
\src{07\_figures.py fig3\_taxonomy\_vs\_aai\_nj}
\tcfig{\figTangle}{%
\textbf{NCBI taxonomy and AAI agree on the large groups.}
Left: NCBI taxonomy cladogram (topology only) over the \nProteomes{} organisms with a
proteome. Right: neighbor-joining tree (Biopython) on $d = -\ln(\mathrm{AAI}/100)$, rooted
on the bacteria, drawn as topology with children rotated to follow the left tree. Lines join
the same organism and take its taxonomic group's color (\texttt{05\_trees.py}).}{fig:tangle}

\clearpage
\src{07\_figures.py fig4\_clade\_tree\_within\_distance; clade\_snp\_distance.csv, clades\_nj.nwk, clade\_distance\_summary.csv}
\tcfig{\figClades}{%
\textbf{Every named \org{S. cerevisiae} clade but one is tighter than the species as a
whole.}
(a) Neighbor-joining tree (Biopython) over \nLeaves{} groups of Peter et al. 2018's
\nIsolates{} isolates: the \nNamedClades{} Peter clades and the \nNoClade{} isolates with no
clade. The distance between two groups is the mean SNP distance over every pair of their
members; midpoint rooted, branch lengths in percent. Bold leaf and black dot: the clade
holding \nearestStrain{}, the isolate with the lowest mean per-ORF nucleotide divergence from
S288C in experiment 018 (\nearestDiv\%); S288C itself is not in Peter's matrix.
(b) SNP distance between every pair of isolates within a group, same rows. \nBelow{} of the
\nNamedWithPairs{} named clades with at least one pair have a median below the median over
all \nAllPairs{} isolate pairs (\medAllPairs\%, dashed gray line); the exception is
\cladeMedMaxName{} (\cladeMedMax\%). Clade medians start at \cladeMedMin\%
(\cladeMedMinName{}), and their median is \medCladeMed\%. The isolates with no clade, a
mixed group, have a median of \noCladeMed\%. Boxes (groups with at least \minBoxPairs{}
pairs): median, quartiles and whiskers at one and a half times the interquartile range,
outliers not drawn; points: every pair of a smaller group. One color for every group, since
color carries no category here. Log axis.}{fig:clades}

\clearpage
\src{07\_figures.py fig5\_ortholog\_fraction\_vs\_aai}
\tcfig{\figOrthologs}{%
\textbf{Shared orthologs fall with AAI.}
Each point is one of the \nPairs{} pairs of the \nProteomes{} proteomes: AAI against the
ortholog fraction (\texttt{03\_mmseqs\_rbh.py}). Color is the pair class, \nClasses{} in
all (n pairs in the legend): within \org{Saccharomyces} \nClsGenus{}, within mammals
\nClsMammals{}, mammal vs chicken \nClsMammalBird{}, within Saccharomycotina
\nClsYcotina{}, within bacteria \nClsBacteria{}, Saccharomycotina vs bacteria
\nClsYeastBacteria{}, Saccharomycotina vs amniote \nClsYeastAnimal{}, amniote vs bacteria
\nClsAnimalBacteria{}, \org{R. toruloides} vs Saccharomycotina \nClsBasidioYeast{}, vs
amniote \nClsBasidioAnimal{} and vs bacteria \nClsBasidioBacteria{}. A within-group class
takes that group's color in Figure~\ref{fig:claim}; a cross-group class takes a palette
color no group uses. Points are drawn largest class first, and the cross-kingdom pairs at
the lowest AAI overlap.}{fig:orthologs}

%% notes-tex/database/044-organism-similarity/sections/headline.tex
%% [[experiments.044-organism-similarity]]
%% https://github.com/Mjvolk3/torchcell/tree/main/notes-tex/database/044-organism-similarity/sections/headline.tex
%% SOURCE: experiments/044-organism-similarity/scripts/08_notes_tex_numbers.py (tables/headline.tex)

\clearpage
\section{Headline numbers\secstatus{todo}}\label{sec:headline}

%% GENERATED FILE -- do not hand-edit.
%% SOURCE: experiments/044-organism-similarity/scripts/08_notes_tex_numbers.py
\begin{table}[H]\centering\footnotesize
\caption[]{Headline numbers. Script column: the numbered script under \texttt{experiments/044-organism-similarity/scripts/} that computed the value; every value is read from \texttt{results/headline\_numbers.json} or \texttt{results/figure\_stats.json}.}
\label{tab:headline}
\begin{tabular}{@{}L{78mm} L{82mm} l@{}}
\toprule
Quantity & Value & Script \\
\midrule
Isolate vs S288C protein identity, median over isolates of per-ORF median & 99.72\% ($n$ = 1,011 isolates, 6,005 ORFs each) & 06 \\
Same, median over isolates of mean over ORFs & 98.04\% (range 95.55 to 99.94\%) & 06 \\
Same, mean over equal-length ORFs only, median over isolates & 99.58\% & 06 \\
One human genome vs GRCh38, peptide-altering sites (literature) & 10,000 to 12,000, i.e. 99.90 to 99.91\% identity & 00 \\
human/chimp AAI, mean / median over RBHs & 98.14 / 99.20\% ($n_\mathrm{RBH}$ 18,461; ortholog fraction 0.920) & 03 \\
human/macaque AAI, mean / median over RBHs & 95.29 / 96.90\% ($n_\mathrm{RBH}$ 18,064; ortholog fraction 0.900) & 03 \\
human/mouse AAI, mean / median over RBHs & 83.70 / 86.50\% ($n_\mathrm{RBH}$ 16,925; ortholog fraction 0.843) & 03 \\
human/chicken AAI, mean / median over RBHs & 72.12 / 73.60\% ($n_\mathrm{RBH}$ 13,170; ortholog fraction 0.731) & 03 \\
S288C/\org{S. paradoxus} AAI, mean / median over RBHs & 92.15 / 93.20\% ($n_\mathrm{RBH}$ 5,367; ortholog fraction 0.969) & 03 \\
S288C/\org{S. eubayanus} AAI, mean / median over RBHs & 83.55 / 84.70\% ($n_\mathrm{RBH}$ 5,136; ortholog fraction 0.955) & 03 \\
SD108/CBS573 AAI, mean / median over RBHs & 99.56 / 99.80\% ($n_\mathrm{RBH}$ 5,037; ortholog fraction 0.980) & 03 \\
S288C/\org{E. coli} AAI & 36.84\% ($n_\mathrm{RBH}$ 391; ortholog fraction 0.091) & 03 \\
S288C/\org{P. kudriavzevii} SD108 AAI & 47.1\% ($n_\mathrm{RBH}$ 3,072 of 5,159 SD108 genes; ortholog fraction 0.595) & 03 \\
Within \org{Saccharomyces} (3 pairs) & AAI median 84.2\%, range 83.5 to 92.1\% & 07 \\
Within mammals (28 pairs) & AAI median 84.3\%, range 83.5 to 98.1\% & 07 \\
Mammal vs chicken (8 pairs) & AAI median 72.1\%, range 71.5 to 72.2\% & 07 \\
Within Saccharomycotina (75 pairs) & AAI median 47.4\%, range 46.0 to 99.6\% & 07 \\
Within bacteria (45 pairs) & AAI median 40.5\%, range 38.8 to 72.7\% & 07 \\
Saccharomycotina vs bacteria (130 pairs) & AAI median 36.5\%, range 35.1 to 38.2\% & 07 \\
Single-copy OG AAI vs RBH AAI, 253 eukaryote pairs & Pearson $r$ 0.998; proxy lower by 1.66 points at the median & 06, 07 \\
Within-clade SNP distance, 28 named clades with a pair & clade medians 0.0015 to 1.03\%; all 510,555 isolate pairs, median 1.02\% & 07 \\
\bottomrule
\end{tabular}
\end{table}


%% notes-tex/database/044-organism-similarity/sections/organisms.tex
%% [[experiments.044-organism-similarity]]
%% https://github.com/Mjvolk3/torchcell/tree/main/notes-tex/database/044-organism-similarity/sections/organisms.tex
%% SOURCE: experiments/044-organism-similarity/scripts/08_notes_tex_numbers.py (tables/organisms.tex)

\clearpage
\section{Organisms\secstatus{todo}}\label{sec:organisms}

%% GENERATED FILE -- do not hand-edit.
%% SOURCE: experiments/044-organism-similarity/scripts/08_notes_tex_numbers.py
\begin{table}[H]\centering\footnotesize
\caption[]{The 34 organisms, 33 with a proteome. Source: \emph{tier} is a torchcell genomes tier set, \emph{NCBI} an assembly fetched with \texttt{datasets}; \emph{Per gene} is the protein count after keeping the longest protein per gene (\texttt{02\_proteome\_prep.py}; \texttt{results/organisms.json}).}
\label{tab:organisms}
\begin{tabular}{@{}l l l l r@{}}
\toprule
Organism & Source & Accession or tier set & Strain & Per gene \\
\midrule
\org{S. cerevisiae} S288C & tier & sgd\_S288C\_R64-4-1\_20230830 & S288C & 6,716 \\
\org{P. kudriavzevii} SD108 & tier & pkudriavzevii\_SD108\_v1 & SD108 & 5,159 \\
\org{P. kudriavzevii} CBS573 & NCBI & GCF\_003054445.1 & CBS573 & 5,139 \\
\org{E. coli} K-12 MG1655 & tier & GCA\_000005845.2 & K-12 MG1655 & 4,290 \\
\org{P. putida} KT2440 & tier & GCA\_000007565.2 & KT2440 & 5,564 \\
\org{C. glutamicum} ATCC 13032 & NCBI & GCF\_000011325.1 & ATCC 13032 & 2,941 \\
\org{B. subtilis} 168 & NCBI & GCF\_000009045.1 & 168 & 4,237 \\
\org{B. licheniformis} DSM 13 & NCBI & GCF\_034478925.1 & ATCC 14580 & 4,169 \\
\org{K. phaffii} GS115 & NCBI & GCF\_000027005.1 & GS115 & 5,040 \\
\org{C. griseus} (hamster) & NCBI & GCF\_003668045.3 & 17A/GY & 21,776 \\
\org{H. sapiens} GRCh38 & NCBI & GCF\_000001405.40 & none & 20,076 \\
\org{L. lactis} IL1403 & NCBI & GCF\_000006865.1 & IL1403 & 2,230 \\
\org{A. limacinum} & absent & GCA\_036850685.1 (no annotation) & none & none \\
\org{Y. lipolytica} CLIB122 & NCBI & GCF\_000002525.2 & CLIB122 & 6,471 \\
\org{R. toruloides} NP11 & NCBI & GCF\_000320785.1 & NP11 & 8,140 \\
\org{C. necator} H16 & NCBI & GCF\_000009285.1 & H16 & 6,618 \\
\org{K. marxianus} DMKU3-1042 & NCBI & GCF\_001417885.1 & DMKU3-1042 & 4,941 \\
\org{H. bluephagenesis} TD01 & NCBI & GCF\_054790765.1 & TD01 & 3,635 \\
\org{O. polymorpha} NCYC 495 & NCBI & GCF\_001664045.1 & NCYC 495 leu1.1 & 5,153 \\
\org{V. natriegens} ATCC 14048 & NCBI & GCF\_001456255.1 & ATCC 14048 & 4,504 \\
\org{M. extorquens} AM1 & NCBI & GCF\_000022685.1 & AM1 & 6,151 \\
\org{L. starkeyi} NRRL Y-11557 & NCBI & GCA\_001661325.1 & NRRL Y-11557 & 8,182 \\
\org{Z. bailii} IST302 & NCBI & GCA\_900162805.1 & IST302 & 5,141 \\
\org{Z. rouxii} CBS 732 & NCBI & GCF\_000026365.1 & CBS 732 & 4,991 \\
\org{L. tetrasporus} NRRL Y-64009 & NCBI & GCF\_029532465.1 & NRRL Y-64009 & 8,004 \\
\org{S. paradoxus} CBS432 & NCBI & GCF\_002079055.1 & CBS432 & 5,538 \\
\org{S. eubayanus} FM1318 & NCBI & GCF\_001298625.1 & FM1318 & 5,379 \\
\org{P. troglodytes} & NCBI & GCF\_028858775.2 & AG18354 & 22,543 \\
\org{G. gorilla} & NCBI & GCF\_029281585.2 & KB3781 & 22,640 \\
\org{P. abelii} & NCBI & GCF\_028885655.2 & AG06213 & 22,203 \\
\org{M. mulatta} & NCBI & GCF\_049350105.2 & MMU2019108-1 & 21,513 \\
\org{M. musculus} C57BL/6J & NCBI & GCF\_000001635.27 & C57BL/6J & 22,198 \\
\org{R. norvegicus} & NCBI & GCF\_036323735.1 & BN/NHsdMcwi & 23,209 \\
\org{G. gallus} & NCBI & GCF\_016699485.2 & bGalGal1 & 18,023 \\
\bottomrule
\end{tabular}
\end{table}


%% notes-tex/database/044-organism-similarity/sections/notdone.tex
%% [[experiments.044-organism-similarity]]
%% https://github.com/Mjvolk3/torchcell/tree/main/notes-tex/database/044-organism-similarity/sections/notdone.tex
%% SOURCE: notes/experiments.044-organism-similarity.md, section "What was not done"

\section{What was not done\secstatus{todo}}\label{sec:notdone}

\begin{itemize}
\item \textbf{BUSCO itself.} The BUSCO-restricted AAI is the single-copy OG proxy;
  compleasm and a current busco were not installable in the experiment's environment.
\item \textbf{\org{A. limacinum}.} No annotated assembly exists at NCBI for
  \org{Aurantiochytrium limacinum} or \org{Schizochytrium}; the organism stays on the
  cladogram (Figure~\ref{fig:heatmaps}) and is excluded from every proteome axis.
\item \textbf{\org{Z. bailii} CLIB 213.} Its GenBank annotation is partial, so the same
  species' strain IST302, with a complete annotation, stands in; a strain substitution, not
  a species one, and the figure labels say IST302.
\item \textbf{CLEAN for EC annotation} is not run; eggNOG-mapper is the annotator.
\item \textbf{Isolates through the mmseqs2 RBH pipeline.} Isolate identity comes from a
  global alignment that charges every gap column, species AAI from local alignments of
  reciprocal best hits; the isolate proteomes were not run through the species pipeline.
\item \textbf{One AAI-matched isolate statistic.} The mean and the median over ORFs are
  both reported and bracket the comparison with species AAI; neither is the same statistic
  as AAI.
\end{itemize}


\end{document}
